% Part: computability
% Chapter: recursive-functions
% Section: partial-functions

\documentclass[../../../include/open-logic-section]{subfiles}

\begin{document}

\olfileid{cmp}{rec}{par}
\olsection{આંશિક પુનરાવર્તી વિધેયો}


પુનરાવર્તી વિધેયોની વ્યાખ્યા તરફ વધતાં પહેલાં નોંધો કે આદિમ
પુનરાવર્તી ન હોય એવાં સંગણનીય વિધેયો અસ્તિત્વમાં છે એવી આપણી
સાબિતી ખરેખર ઘણું વધુ બતાવે છે. દલીલ સરળ હતી: આપણે માત્ર એ
હકીકત વાપરી કે વિધેયો $f_0,f_1,\dots$ને એવી રીતે પરિગણિત કરી
શકાય કે $x$ અને $y$ના વિધેય તરીકે $f_x(y)$ સંગણનીય હોય. તેથી
આ દલીલ \emph{આ રીતે પરિગણિત કરી શકાય એવા વિધેયોના કોઈ પણ વર્ગ}
માટે લાગુ પડે છે. હવે મુશ્કેલી એ છે કે આપણે સંગણનીય વિધેયોનું
સ્પષ્ટ વર્ણન કરવા માગીએ છીએ, પણ સંગણનીય વિધેયોના સંગ્રહનું કોઈ
પણ આવું સ્પષ્ટ વર્ણન સંપૂર્ણ હોઈ શકે નહીં!

ઉકેલ એ છે કે \emph{આંશિક} વિધેયોને પણ સામેલ કરીએ. આપણે જોઈશું કે
આંશિક સંગણનીય વિધેયોની પરિગણના \emph{શક્ય} છે.
\iftag{TMs}{ ખરેખર, ટ્યુરિંગ મશીનોને વ્યવસ્થિત રીતે પરિગણિત
  કરી શકાય છે, તેથી આ વાત લગભગ આપણે પહેલેથી જાણીએ છીએ.}{}
પછી આપણે વિકર્ણ-દલીલ તરફ પાછા ફરીશું અને આંશિક વિધેયોનો
સમાવેશ કરતાં એ દલીલ કેમ ચાલતી નથી તે તપાસીશું.

હવે પ્રશ્ન એ છે કે બધા આંશિક પુનરાવર્તી વિધેયો મેળવવા માટે
આદિમ પુનરાવર્તી વિધેયોમાં શું ઉમેરવું જોઈએ? બે બાબતો કરવી પડશે:
\begin{enumerate}
\item આદિમ પુનરાવર્તી વિધેયોની વ્યાખ્યા બદલીને તેમાં આંશિક
  વિધેયોને પણ આવવા દેવા.
\item વ્યાખ્યામાં કંઈક \emph{ઉમેરવું}, જેથી નવા આંશિક વિધેયોનો
  સમાવેશ થાય.
\end{enumerate}

પહેલું સહેલું છે. પહેલાંની જેમ શૂન્ય વિધેય, ઉત્તરગામી વિધેય અને
પ્રક્ષેપોથી શરૂ કરીને સંયોજન અને આદિમ પુનરાવર્તન હેઠળ બંધ કરીશું.
ફરક માત્ર એટલો કે વ્યાખ્યામાંના કેટલાક પદોનું મૂલ્ય ન હોય એવી
શક્યતા સમાવવા સંયોજન અને આદિમ પુનરાવર્તનની વ્યાખ્યાઓ બદલવી પડશે.
$f$ અને $g$ આંશિક વિધેયો હોય, તો $f(x) \fdefined$થી જણાવીશું કે
$f$ એ $x$ પર વ્યાખ્યાયિત છે, એટલે $x$ એ $f$ના પ્રદેશમાં છે;
અને $f(x)
\fundefined$થી તેનો વિરોધી અર્થ, એટલે $f$ એ~$x$ પર
વ્યાખ્યાયિત નથી, જણાવીશું. $f(x) \simeq g(x)$થી જણાવીશું કે
$f(x)$ અને $g(x)$ બંને અવ્યાખ્યાયિત છે, અથવા બંને વ્યાખ્યાયિત
અને સમાન છે. વધુ જટિલ પદો માટે પણ આ સંકેતનો વાપરીશું. એક
રૂઢિ અપનાવીએ: $h$ અને $g_0$, \dots,~$g_k$ બધાં આંશિક
વિધેયો હોય, તો
\[
h(g_0(\vec x),\dots,g_k(\vec x))
\]
ત્યારે અને માત્ર ત્યારે જ વ્યાખ્યાયિત છે જ્યારે દરેક $g_i$ એ
$\vec x$ પર વ્યાખ્યાયિત હોય અને $h$ એ
$g_0(\vec x)$, \dots,~$g_k(\vec x)$ પર વ્યાખ્યાયિત હોય. આ
સમજ સાથે આંશિક વિધેયો માટે સંયોજન અને આદિમ પુનરાવર્તનની
વ્યાખ્યાઓ અગાઉ જેવી જ છે; માત્ર ``$=$''ની જગ્યાએ
``$\simeq$'' મૂકવું પડે છે.

આદિમ પુનરાવર્તી વિધેયોની વ્યાખ્યામાં આંશિક વિધેયો મેળવવા જે
ઉમેરીશું તે \emph{અનિયત મર્યાદા સુધીની શોધક્રિયા} છે. પ્રાકૃતિક
સંખ્યાઓ પર $f(x,\vec z)$ કોઈ પણ આંશિક વિધેય હોય, તો
$\mu x \; f(x,\vec z)$ને આ રીતે વ્યાખ્યાયિત કરીએ:
\begin{quote}
  જો એવો $x$ અસ્તિત્વમાં હોય કે $f(0,\vec z), f(1,\vec z),
  \dots, f(x,\vec
  z)$ બધાં વ્યાખ્યાયિત હોય અને $f(x,\vec z) = 0$ હોય, તો એવા
  $x$માંથી સૌથી નાનો.
\end{quote}
અન્યથા $\mu x \; f(x,\vec z)$ અવ્યાખ્યાયિત છે એમ સમજવું.
આથી $\mu x \; f(x,\vec z)$ એકમાત્ર રીતે વ્યાખ્યાયિત થાય છે.

\begin{explain}
નોંધો કે આપણી વ્યાખ્યા\iftag{TMs}{ ટ્યુરિંગ મશીનો, અથવા}{}
અલ્ગોરિધમો કે સંગણનના કોઈ ખાસ નમૂનાનો ઉલ્લેખ કરતી નથી. પરંતુ
સંયોજન અને આદિમ પુનરાવર્તનની જેમ અનિયત મર્યાદા સુધીની શોધ પાછળ
ક્રિયાત્મક, સંગણનાત્મક સમજ છે. આંશિક વિધેયની સંગણનીયતામાં જે
આર્ગ્યુમેન્ટો પર વિધેય અવ્યાખ્યાયિત છે તે એવા ઇનપુટને અનુરૂપ છે
જેના પર ગણતરી અટકતી નથી. $\mu x \;
f(x,\vec z)$ ગણવાની રીત આ છે: $f(0,\vec z), f(1,\vec z),
f(2,\vec z)$ વગેરેની ગણતરી 0 મૂલ્ય મળે ત્યાં સુધી કરો.
વચલી કોઈ ગણતરી ન અટકે, તો $\mu x \; f(x,\vec z)$ની ગણતરી પણ
નહીં અટકે.
\end{explain}

$R(x,\vec z)$ કોઈ પણ સંબંધ હોય, તો $\mu x \; R(x,\vec z)$ને
$\mu x \; (1 \tsub \Char{R}(x,\vec z))$ તરીકે વ્યાખ્યાયિત
કરીએ છીએ. બીજા શબ્દોમાં, $\mu x \; R(x,\vec z)$ એ $R(x,\vec z)$
લાગુ પડે એવું $x$નું સૌથી નાનું મૂલ્ય આપે છે. તેથી $f(x,\vec z)$
સંપૂર્ણ વિધેય હોય, તો $\mu x \; f(x,\vec
z)$ એ $\mu x \; (f(x,\vec z) = 0)$ જેવું જ છે. જોકે આપણી
મૂળ વ્યાખ્યા વધુ સામાન્ય છે: તેમાં $f(x,\vec z)$ દરેક જગ્યાએ
વ્યાખ્યાયિત ન હોય એવી છૂટ છે; તેની સામે સંબંધનું લક્ષણ વિધેય
હંમેશાં સંપૂર્ણ હોય છે.

\begin{defn}
\emph{આંશિક પુનરાવર્તી વિધેયો}નો ગણ એટલે પ્રાકૃતિક સંખ્યાઓમાંથી
પ્રાકૃતિક સંખ્યાઓમાં જતાં (જુદી જુદી આર્ગ્યુમેન્ટ-સંખ્યા ધરાવતાં)
આંશિક વિધેયોનો સૌથી નાનો એવો ગણ, જેમાં શૂન્ય વિધેય,
ઉત્તરગામી વિધેય અને પ્રક્ષેપો છે તથા જે સંયોજન, આદિમ પુનરાવર્તન
અને અનિયત મર્યાદા સુધીની શોધ હેઠળ બંધ છે.
\end{defn}

અલબત્ત, કેટલાંક આંશિક પુનરાવર્તી વિધેયો સંપૂર્ણ પણ નીવડે છે,
એટલે દરેક આર્ગ્યુમેન્ટ માટે વ્યાખ્યાયિત હોય છે.

\begin{defn}
\ollabel{defn:recursive-fn}
\emph{પુનરાવર્તી વિધેયો}નો ગણ એટલે સંપૂર્ણ હોય એવા આંશિક
પુનરાવર્તી વિધેયોનો ગણ.
\end{defn}

પુનરાવર્તી વિધેયને ક્યારેક ``સંપૂર્ણ પુનરાવર્તી'' પણ કહે છે,
જેથી તે દરેક જગ્યાએ વ્યાખ્યાયિત છે તે પર ભાર પડે.

\end{document}
